\chapter{Assembly Meets C: One Game, Two Altitudes}
\label{ch:join}

This short chapter is the book's reunion. If you have read both tracks, it will
crystallise something you have half-seen already; if you have read only one, it is
an invitation to glance over the wall at the other side. The claim is simple and,
once shown, hard to unsee: the assembly game and the C game are the same game,
viewed at two altitudes, and the compiler is the translator between them
\citep{prg32teaching}.

\section{The same logic, two languages}

Consider the single most common operation in any of our games: move an object left
when the left button is pressed. In C it is three lines you could read aloud as
English.

\begin{lstlisting}[style=c,caption={Move left, in C. Read input, test the bit, change the position.}]
uint32_t input = prg32_input_read();
if (input & PRG32_BTN_LEFT) {
    paddle_x -= 3;
}
\end{lstlisting}

In assembly it is the same idea spelled out in the machine's own terms: call,
copy the result, mask the bit, branch if clear, subtract, and---because
\code{paddle\_x} lives in memory---load before and store after.

\begin{lstlisting}[style=asm,caption={Move left, in assembly. The same logic, with every step the C compiler would have generated made explicit.}]
    call prg32_input_read     # input bitmask in a0
    mv   t0, a0               # keep our copy
    la   t1, paddle_x
    lw   t2, 0(t1)            # load paddle_x from memory
    andi t3, t0, 1            # isolate the LEFT bit
    beqz t3, 1f               # if LEFT is clear, skip
    addi t2, t2, -3           # paddle_x -= 3
1:  sw   t2, 0(t1)            # store paddle_x back
\end{lstlisting}

Line up the two listings and the correspondence is exact. The C \code{if (input \&
PRG32\_BTN\_LEFT)} becomes \code{andi} followed by \code{beqz}. The C \code{paddle\_x
-= 3} becomes a load, an \code{addi}, and a store, because a C variable in memory
must be brought into a register to be changed and written back---the
load--modify--store cycle of Chapter~\ref{ch:asm1}. Nothing in the assembly is
arbitrary; every instruction is doing work the C statement implied.

\begin{prgconcept}
A compiler translates each C construct into a fixed pattern of machine
instructions: a bitwise test becomes a mask and a branch, an assignment to a
variable in memory becomes load--modify--store, a function call becomes the
calling convention. C is not magic that the machine somehow understands; it is a
notation that a compiler mechanically turns into exactly the assembly you would
have written by hand.
\end{prgconcept}

\section{What C hides, and what it does not}

The C version is shorter not because it does less but because the compiler writes
the bookkeeping for you. It chooses the registers, emits the loads and stores,
manages the stack frame, and saves \reg{ra} around calls---all the discipline you
practised by hand in Part~II. This is C's bargain: you express intent, the
compiler handles the mechanism.

What C does \emph{not} hide, on this platform, is the machine itself. The same
display, the same input bitmask, the same RGB565 colours, the same three-function
contract, the same two targets. A C variable is still a box in memory; a collision
is still a comparison; a frame is still a turn around a loop. The abstractions of C
sit directly on the concepts of Part~I, which is why a student who has seen both
tracks understands C more deeply than one who has only ever used it: they know what
the compiler is doing on their behalf.

\section{A side-by-side map}

The platform's teaching guide recommends exactly this comparison as a classroom
activity, using the platformer example, whose assembly and C versions implement
the same game \citep{prg32teaching}. Table~\ref{tab:sidebyside} lays the two
languages alongside each other for the concepts you have met.

\begin{center}
\begin{tabularx}{\linewidth}{@{}lXX@{}}
\toprule
Concept & In assembly (Part~II) & In C (Part~III) \\
\midrule
Persistent state & labelled words in \code{.data} & file-scope variables \\
A grouped object & a block of memory at fixed offsets & a \code{struct} \\
Many objects & an array of words & an array \\
Read input & \code{call}, then \code{andi}/\code{beqz} & \code{prg32\_input\_read} and \code{if (\dots \& \dots)} \\
Change a variable & load, modify, store & \code{x += \dots} \\
A decision & a conditional branch & an \code{if} \\
Call the framework & manage \reg{ra} and the stack & write the call \\
Move an actor & \code{call prg32\_platform\_actor\_step} & the same call, by name \\
\bottomrule
\end{tabularx}
\captionof{table}{The same ideas in both languages. The right column is the left
column with the bookkeeping delegated to the compiler
\citep{prg32teaching}.}
\label{tab:sidebyside}
\end{center}

\begin{prgaside}
That the framework call \code{prg32\_platform\_actor\_step} appears identically in
both columns is the whole point of the platform's design. The framework is written
in C and exposes a standard \riscv calling convention, so an assembly game and a C
game invoke it the same way. The boundary between ``your code'' and ``the
framework'' is the same boundary in both languages, which is why the two tracks can
share a single runtime, a single set of examples, and a single pair of targets
\citep{prg32abi,prg32teaching}.
\end{prgaside}

\section{Where to go next}

You have now written native \riscv games---in assembly, in C, or in both---that run
on a real open-architecture processor and on its emulator. The natural next steps
are to read the richer examples you have not yet opened: Space Invaders for many
moving enemies, Tetris for a falling-piece grid, the raycaster for fixed-point
pseudo-3-D, and Wing Commander for two layered playfields \citep{prg32repo}. Each is
a variation on patterns you already know, and each is written to be read.

Beyond the games, the appendices that follow are a working reference for everything
you have used and several things you have not: the complete framework interface,
the workflow for measuring a game's performance like a scientist, full annotated
source for the examples, the use of GitHub to manage and submit your work, the
construction of the physical board, and the setup of your development environment
on whichever operating system you use.

\begin{prgcheck}
You should now be able to look at a short C statement and describe, in words, the
assembly the compiler will produce for it---and to look at a fragment of \riscv
assembly and recognise the C construct it implements. That two-way fluency, between
the notation and the machine, is the deepest thing a first year of computing can
give you, and it is what this book set out to teach.
\end{prgcheck}
